\chapter{Cervical Cerclage}

\section*{Overview}
Cervical cerclage is a procedure that places a suture around the cervix during pregnancy to hold it closed and prevent preterm birth in women with cervical insufficiency. The exact approach, setting, and preparation depend on the indication, anatomy, and the patient's health.

\section*{Alternative Names}
Cervical stitch; cervical suture; cervical incompetence treatment

\section*{Principle}
A non-absorbable suture encircles the cervix at the level of the internal os, providing mechanical support that keeps the cervix closed against progressive dilation during pregnancy.

\section*{History}
Cervical cerclage was introduced in the 1950s (Shirodkar 1955, McDonald 1957) and remains the standard treatment for cervical insufficiency.

\section*{Why It Is Done}
Performed for a history of painless second-trimester loss, a short cervix on ultrasound, or premature cervical dilation in the current pregnancy.

\section*{Is This a Primary or Secondary Test or Procedure}
Primary preventive procedure for cervical insufficiency.

\section*{How It Is Performed}
Performed under regional or general anesthesia. Through the vagina, the cervix is exposed and a suture is placed (McDonald purse-string or Shirodkar technique) and tied to close the internal os.

\section*{Preparation}
Ultrasound evaluation of fetal anatomy and cervical length, and exclusion of infection or preterm labor.

\section*{Contraindications and Precautions}
Chorioamnionitis, preterm premature rupture of membranes, fetal anomaly, and active labor are contraindications.

\section*{Risks}
Infection, membrane rupture, bleeding, cervical injury, and preterm labor.

\section*{Aftercare and Recovery}
The patient is monitored and activity may be limited; the stitch is removed at 36-37 weeks or at delivery.

\section*{Outcome and Follow-up}
Serial ultrasound checks cervical length; the stitch is assessed for displacement.

\section*{Expected Results}
A closed cervix and a pregnancy carried to term.

\section*{Follow-up}
Serial cervical length ultrasound; removal of the stitch near term.

\section*{When to Seek Medical Care}
Follow the procedure team's specific instructions. Seek urgent medical assessment for severe or worsening pain, uncontrolled bleeding, fever or signs of infection, trouble breathing, chest pain, fainting, new weakness or confusion, inability to urinate, or any rapidly worsening symptom. The appropriate warning signs depend on the procedure and the patient's medical condition.

\section*{References}
\begin{enumerate}
  \item MedlinePlus. Cervical cerclage. U.S. National Library of Medicine; 2025.
  \item Current Medical Diagnosis and Treatment. McGraw-Hill; 2025.
\end{enumerate}

\clearpage
